% page 102 du fac-simile (image p-101 du scan)
% bloc 157.5 x 250.6 mm ; marge gauche 8.0 mm
\pgc{-12.700mm}{0.000mm}{
\l{}
\l{\cel{8}94}
\l{}
\l{}
\l{\sou{citadelo}. Kastelo fuorta qua protektas urbo. - DEFIRS.}
\l{}
\l{\sou{citaro}. (\sou{Muziko}). Sorto di muzik-instrumento kordoza, quan}
\l{\cel{3}uzis la antiqui. - DEFIRS.}
\l{}
\l{\sou{citizo}. (\sou{bot.}) Genero de planti di qua la tipo esas la Alpo-}
\l{\cel{3}citizo o \textquotedbl{}falsa ebeniero\textquotedbl{}. - L.\cel{1}\sou{cytisus}.\cel{2}DEFIRS.}
\l{}
\l{\sou{citodo}. (\sou{biol.)}\cel{2}Maso protoplasma nuda, sen-nuklea (ex.: la}
\l{\cel{3}monero di Haeckel).}
\l{}
\l{\sou{citologio}. (\sou{biol.}) Parto di histologio di qua la studio-temo}
\l{\cel{3}esas la celuli, de la vidopunti konstituceso, morfologio}
\l{\cel{3}ed evoluciono. - DEFIRS.}
\l{}
\l{\sou{citoplasmo.}\cel{1}(\sou{biol.}) Protoplasmo, situita éxter la nukleo di}
\l{\cel{3}la celulo vivanta. - DEFIRS.}
\l{}
\l{\sou{citotropismo}. (\sou{biol.}) Atrakto di celulo ad altra celuli viv-}
\l{\cel{3}anta, od ad ula substanci kemiala sen-viva. - DEFIRS.}
\l{}
\l{\sou{citrono}. (\sou{bot.}) Frukto flava klare, saporanta acide, de ar-}
\l{\cel{3}boro di la genero \textquotedbl{}oranjo\textquotedbl{}. - L.\cel{1}\sou{citrus}. - DEF}
\l{}
\l{\sou{civeto}. (\sou{zool.}) Mamifero karnivora, analoga a la martro, qua}
\l{\cel{3}sekrecas quaza mosko. - L.\cel{1}\sou{viverra zibetha}. - DEFIRS.}
\l{}
\l{\sou{civila}. Qua relatas la civitani. (Antonimo di armeala, reli-}
\l{\cel{3}giala, kriminala, politikala). -\cel{1}\sou{Estado civila}\cel{1}: ensemblo}
\l{\cel{3}de la dokumenti relativa a la estado sociala di singla}
\l{\cel{3}civitano : naskala akto, mariajala akto, e c. - DEFIRS.}
\l{}
\l{\sou{civilizar}.\cel{2}(\sou{trans.}) Pasigar de la estado primitiva, naturala,}
\l{\cel{3}ad estado plu progresoza pri la kultiveso etikala, morala,}
\l{\cel{3}pensala e sociala. - DEFIRS.}
\l{}
\l{\sou{civismo}. Devoteso di la civitano a la Stato. - DEFIS.}
\l{}
\l{\sou{civito}. I. Urbo, konsiderata kom korpo politikala. - II. La}
\l{\cel{3}korpo ek la civitanaro. - IS.}
\l{}
\l{\sou{cizo}. Instrumento qua konsistasye du lami di qui la bordi}
\l{\cel{3}tranchanta interkrucumas, ed uzata por tranchar la kozi}
\l{\cel{3}dina (stofo, papero, filo, e c.) - EFI.}
\l{}
\l{\sou{cizelar}. (\sou{trans.}) Laborar metalajo (bronzo, oro, arjento,}
\l{\cel{3}e c.) per instrumento ek stalo, longa, plata, tranchiva}
\l{\cel{3}ye un extremajo e quan onu uzas por entaliar e nochizar.}
\l{\cel{3}- DEFIRS.}
\l{}
\l{co. Ca kozo.}
\l{}
\l{}
\l{\cel{30}----}
}
